\begin{abstract}
Satellite Internet of Things services reach remote rivers but sell only a few short, one-way messages per day. This study developed a framework that delivers continuous water-quality records within the accuracy that the U.S. Geological Survey (USGS) requires of an excellent record and lets the server certify, from the received bytes alone, which readings meet that tolerance. A node sends the fewest values that keep every reading within tolerance, in 23-byte messages protected by an erasure code sized to each day's record; this choice is optimal among linear designs without knowledge of the loss process. After development on 39 USGS sites, the frozen framework was confirmed on 36 sealed sites in five other regions. A four-parameter day needed a median of two messages. With six messages per day and 30\% simulated loss, the server fully certified \gCodeIidConf{} of days, ahead of repetition, a fixed code, and satellite-tag practice; twelve messages per day certified \gTwelveIidConf{}. Certified records never gave a wrong certain decision and caught conductance pulses that periodic telemetry missed, but did not resolve features finer than the tolerance. Six of eight confirmatory hypotheses held.
\end{abstract}

\begin{keyword}
water-quality monitoring \sep satellite Internet of Things \sep error-bounded records \sep erasure coding \sep data certification \sep reproducible evaluation
\end{keyword}
